% !TeX program = pdflatex
% Feuille TD 2 - Licence 1 tronc commun - Analyse 2MTH1 0101
\documentclass[11pt,a4paper,oneside]{article}
\usepackage[T1]{fontenc}
\usepackage[utf8]{inputenc}
\usepackage[main=french,provide=*]{babel}
\usepackage{lmodern,microtype,amsmath,amssymb,mathtools}
\usepackage{geometry}
\geometry{left=2.4cm,right=2.3cm,top=2.55cm,bottom=2.45cm,headheight=24pt}
\usepackage{enumitem,booktabs,tabularx,needspace}
\usepackage[most]{tcolorbox}
\usepackage{fancyhdr,lastpage,titlesec,hyperref}
\definecolor{Marine}{HTML}{10283F}
\definecolor{Bleu}{HTML}{176D8E}
\definecolor{Sarcelle}{HTML}{118F86}
\definecolor{Or}{HTML}{EDAA47}
\definecolor{GrisTexte}{HTML}{54667A}
\definecolor{BleuPale}{HTML}{EFF5F8}
\hypersetup{colorlinks=true,linkcolor=Bleu,urlcolor=Sarcelle,pdftitle={Analyse L1 - TD 2 - Calcul differentiel},pdfauthor={M. Bakary TRAORE}}
\setlength{\parindent}{0pt}
\setlength{\parskip}{5pt plus1pt minus1pt}
\setlist[enumerate]{leftmargin=1.55em,itemsep=3pt,topsep=3pt}
\titleformat{\section}{\Large\bfseries\color{Marine}}{}{0pt}{}[\vspace{1mm}{\color{Sarcelle}\titlerule[.8pt]}]
\titleformat{\subsection}{\large\bfseries\color{Bleu}}{}{0pt}{}
\titlespacing*{\section}{0pt}{19pt}{8pt}
\titlespacing*{\subsection}{0pt}{12pt}{4pt}
\pagestyle{fancy}\fancyhf{}
\fancyhead[L]{\small\color{Marine}\textbf{UNIVERSITÉ DE N'ZÉRÉKORÉ}}
\fancyhead[R]{\small\color{GrisTexte}Analyse L1 \textbullet\ 2MTH1 0101}
\fancyfoot[L]{\scriptsize\color{GrisTexte}TD 2 \textbullet\ Calcul différentiel}
\fancyfoot[R]{\scriptsize\color{GrisTexte}\thepage\,/\,\pageref{LastPage}}
\newtcolorbox{objectifbox}{enhanced,breakable,colback=BleuPale,colframe=Bleu,boxrule=.6pt,arc=1mm}
\newcommand{\R}{\mathbb R}
\newcommand{\e}{\mathrm e}
\begin{document}
\begin{tcolorbox}[enhanced,colback=Bleu,colframe=Bleu,boxrule=0pt,sharp corners,left=5mm,right=5mm,top=5mm,bottom=5mm]
\begin{center}
{\color{white}\large\bfseries UNIVERSITÉ DE N'ZÉRÉKORÉ\par}\vspace{.14cm}
{\color{white!85}\normalsize LICENCE 1 -- TRONC COMMUN \quad / \quad 2MTH1 0101\par}\vspace{.30cm}
{\color{white}\LARGE\bfseries FEUILLE DE TRAVAUX DIRIGÉS (TD 2)\par}\vspace{.14cm}
{\color{white!95}\large CHAPITRE 2 -- Calcul différentiel\par}\vspace{.14cm}
{\color{white!80}\small M. Bakary TRAORE \quad | \quad 2026--2027\par}
\end{center}
\end{tcolorbox}
\begin{objectifbox}
\textbf{Durée :} 2 h dirigées (1 h par grandes parties, exercices complémentaires à terminer en autonomie).\\
\textbf{Consignes :} indiquer les domaines et les unités, justifier l'usage des théorèmes, détailler les calculs. Les exercices $\star$ sont des prolongements.\\
\textbf{Filières :} Météorologie, Hydrologie, Génie de l'environnement et Gestion des Ressources Naturelles.
\end{objectifbox}

\section*{Partie A -- Dérivation et dérivées successives (35 min)}
\Needspace{5\baselineskip}
\subsection*{Exercice 1 -- Règles de dérivation}
Déterminer le domaine et calculer la dérivée de chaque fonction :
\begin{enumerate}[label=\alph*)]
\item $f(x)=(x^2+1)\e^{-x}$ ;
\item $g(x)=\ln(1+x^2)$ ;
\item $h(x)=\dfrac{\sqrt{2x+3}}{x-1}$.
\end{enumerate}

\Needspace{5\baselineskip}
\subsection*{Exercice 2 -- Dérivée par définition et tangente}
Soit $f(x)=x^2$. En utilisant la définition de $f'(a)$, calculer $f'(2)$, puis écrire l'équation de la tangente en $(2,f(2))$. La fonction $u(x)=|x|$ est-elle dérivable en $x=0$ ? Justifier.

\Needspace{5\baselineskip}
\subsection*{Exercice 3 -- Dérivées successives}
Pour $f(x)=\e^{2x}$, calculer $f'(x)$, $f''(x)$ et $f'''(x)$ ; conjecturer et démontrer par récurrence une formule pour $f^{(n)}(x)$, $n\in\mathbb N$.

\section*{Partie B -- Théorèmes et extrema (45 min)}
\Needspace{5\baselineskip}
\subsection*{Exercice 4 -- Théorème de Rolle}
Vérifier que $f(x)=x^2-4x+3$ satisfait le théorème de Rolle sur $[1,3]$, puis déterminer le point $c$ annoncé par ce théorème.

\Needspace{5\baselineskip}
\subsection*{Exercice 5 -- Théorème des accroissements finis}
Appliquer le TAF à $f(x)=\ln x$ sur $[1,\e]$. Déterminer exactement le nombre $c\in]1,\e[$. Expliquer ce que représente la pente moyenne.

\Needspace{5\baselineskip}
\subsection*{Exercice 6 -- Étude complète d'une fonction}
Soit $f(x)=x^3-3x$ sur $[-2,2]$.
\begin{enumerate}[label=\alph*)]
\item Étudier le signe de $f'$ et les variations ; trouver les extrema locaux.
\item Trouver les minima et maxima \emph{globaux} sur $[-2,2]$.
\item Calculer $f''$ et déterminer les intervalles de concavité et de convexité et le point d'inflexion.
\end{enumerate}

\section*{Partie C -- DL et approximations (20 min)}
\Needspace{5\baselineskip}
\subsection*{Exercice 7 -- Développements limités}
\begin{enumerate}[label=\alph*)]
\item Calculer $\displaystyle\lim_{x\to0}\frac{\e^x-1-x}{x^2}$ et $\displaystyle\lim_{x\to0}\frac{\ln(1+x)-x}{x^2}$.
\item Approcher $\e^{-0.1}$ par le DL d'ordre 2. Donner une majoration de l'erreur à l'aide de Taylor--Lagrange, sachant que $|f^{(3)}(x)|\le\e^{0.1}$ sur $[-0.1,0]$ pour $f(x)=\e^x$.
\end{enumerate}

\section*{Partie D -- Applications dans les quatre filières (20 min)}
\Needspace{5\baselineskip}
\subsection*{Exercice 8 -- Météorologie et hydrologie}
\begin{enumerate}[label=\alph*)]
\item $T(t)=25+3\sin(\pi t/12)$, $T$ en $^\circ$C et $t$ en heures. Calculer $T'(0)$ et interpréter son unité.
\item $Q(t)=5+2t\e^{-t}$, $Q$ en $\mathrm{m^3/s}$ et $t\ge0$ en heures. Déterminer l'heure où le débit est maximal.
\end{enumerate}

\Needspace{5\baselineskip}
\subsection*{Exercice 9 -- Génie de l'environnement}
La concentration d'un polluant suit $C(t)=10\e^{-0.2t}$, $t$ en heures et $C$ en mg/L. Calculer $C'(t)$ et $C'(0)$. Estimer $C(0.5)$ par le DL d'ordre 2 en $0$ et préciser le sens physique du signe de $C'$.

\Needspace{5\baselineskip}
\subsection*{Exercice 10 -- Gestion des Ressources Naturelles $\star$}
La biomasse suit $B(t)=100-20\e^{-0.1t}$, $t$ en années et $B$ en tonnes. Étudier $B'(t)$ et $B''(t)$ pour $t\ge0$. La biomasse tend-elle vers une limite ? Interpréter le résultat.

\vspace{.3cm}\noindent{\small\color{GrisTexte}\textbf{Suite du travail :} consulter le corrigé détaillé \emph{Solutions TD 2} après la séance, puis réaliser le TP 1 (représentation graphique et dérivée numérique).}
\end{document}
